% Adapted from ne630_problems/lesson_17/statement.tex at commit e1fba66.
% Corrected the inconsistent fuel number density in the legacy substitution.

\noindent{\bf Problem 1 [FNRP 4.6].}

A PWR has 3 atom-\% enriched UO$_2$ fuel pins with $D=1.0~\mathrm{cm}$,
$\rho_f=11.0~\mathrm{g\,cm^{-3}}$, and $V_m/V_f=2$. At room temperature,
calculate $\eta_T$, $p$, $f$, and $k_\infty$ for $\epsilon=1.24$,
$\varsigma=1.16$, and an effective resonance diameter $1.1D$.

\begin{adaptedsolution}
The UO$_2$ molecule density and water molecule density are
\begin{align*}
N_f&=\frac{11N_A}{0.03(235)+0.97(238)+2(16)}
=2.454\times10^{22}~\mathrm{cm^{-3}},\\
N_m&=\frac{(1.00)N_A}{18}=3.346\times10^{22}~\mathrm{cm^{-3}}.
\end{align*}
Using the thermal data from FNRP Problem 3.4,
\[
\eta_T=
\frac{2.43(505)/591}
{1+\dfrac{0.97(2.42)}{0.03(591)}}
=\boxed{1.8336}.
\]
With $\sigma_{a,m}=2(0.295)+1.69\times10^{-4}$ b and
$\sigma_{a,f}=0.97(2.42)+0.03(591)$ b,
\[
f=\left[1+\varsigma\frac{V_mN_m}{V_fN_f}
\frac{\sigma_{a,m}}{\sigma_{a,f}}\right]^{-1}
=\boxed{0.9149}.
\]
For the epithermal water scattering values $\sigma_s^H=19.508$ b and
$\sigma_s^O=3.85$ b,
\[
\xi=\frac{2(1)(19.508)+(0.1199)(3.85)}{2(19.508)+3.85}
=0.9210.
\]
The Table 4.3 resonance-integral approximation gives
\[
I=4.45+26.6\sqrt{\frac{4}{\rho_f(1.1D)}}=19.7439~\mathrm{b},
\]
and Eq.~(4.40) then gives
\[
p=\exp\left[-\frac{0.97}{(V_mN_m)/(V_fN_f)}
\frac{I}{\xi[2\sigma_s^H+\sigma_s^O]}\right]
=\boxed{0.8370}.
\]
Finally,
\[
k_\infty=\eta_T f p\epsilon
=(1.8336)(0.9149)(0.8370)(1.24)
=\boxed{1.741}.
\]
The result is high compared with a representative operating PWR because this
room-temperature estimate uses simplified thermal averaging and resonance
self-shielding.
\end{adaptedsolution}

\newpage
\noindent{\bf Problem 2 [FNRP 4.8].}

A uranium lattice's heavy-water moderator is replaced with light water.

\begin{adaptedsolution}
\begin{enumerate}
\item The resonance escape probability $p$ generally \emph{increases} because
light water has greater slowing-down power, $\xi\Sigma_s$, so neutrons spend
less slowing-down time in the resonance-energy range.
\item The thermal utilization $f$ \emph{decreases} because ${}^{1}$H absorbs
substantially more thermal neutrons than ${}^{2}$H.
\item These effects oppose one another, so the sign cannot be established
from the four-factor formula alone without magnitudes. For a lattice designed
for heavy water, a direct replacement by light water will generally reduce
$k_\infty$ because the moderator-absorption penalty is large.
\end{enumerate}
\end{adaptedsolution}

\newpage
\noindent{\bf Problem 3 [FNRP 4.9].}

The coolant-to-fuel volume ratio is increased in a PWR.

\begin{adaptedsolution}
\begin{enumerate}
\item The basic Chapter 4 formulas do not uniquely quantify $\epsilon$.
Physically, added moderation softens the spectrum and usually \emph{decreases}
the fast-fission factor slightly because fewer fast neutrons cause
${}^{238}$U fission.
\item The resonance escape probability $p$ \emph{increases}: increasing
$V_m/V_f$ increases the slowing-down power available per fuel atom and makes
capture during passage through the ${}^{238}$U resonances less likely.
\end{enumerate}
This does not by itself determine the behavior of $k_\infty$, since the
thermal-utilization penalty also changes with moderator amount.
\end{adaptedsolution}
